% !TeX program = pdflatex
%
% "A manufactured Apery apparatus for a cusp-form L-value."
% Working paper of River's odd-zeta programme.
%
% Provenance: work/CUSPIDAL_COMPANION.md (level-6 experiment, Sol Project D),
% work/D1_CUSPIDAL_APPARATUS.md (level-12 construction, delta control,
% uniqueness corollary, proof notes section 8, irrationality race section 9),
% work/PHI_SOURCE_LEDGER.md (the Eisenstein/cuspidal dichotomy of sources),
% scripts work/z5eps/eps60--eps62b.  Companion paper (the proved base):
% papers_out/modular_anchors/main.tex.
%
% Evidence labels are load-bearing and are never silently upgraded.
%
\documentclass[11pt,reqno]{amsart}

\usepackage[T1]{fontenc}
\usepackage[utf8]{inputenc}
\usepackage{amsmath,amssymb,amsthm}
\usepackage{mathtools}
\usepackage{booktabs}
\usepackage{enumitem}
\usepackage[margin=1.05in]{geometry}
\usepackage{microtype}
\usepackage[hidelinks]{hyperref}

\theoremstyle{plain}
\newtheorem{theorem}{Theorem}[section]
\newtheorem{proposition}[theorem]{Proposition}
\newtheorem{lemma}[theorem]{Lemma}
\newtheorem{corollary}[theorem]{Corollary}
\newtheorem{conjecture}[theorem]{Conjecture}

\theoremstyle{definition}
\newtheorem{definition}[theorem]{Definition}
\newtheorem{verification}[theorem]{Computer-assisted finite verification}
\newtheorem{construction}[theorem]{Construction}
\newtheorem{problem}[theorem]{Problem}

\theoremstyle{remark}
\newtheorem{remark}[theorem]{Remark}
\newtheorem{observation}[theorem]{Observation}

\newcommand{\Z}{\mathbb{Z}}
\newcommand{\Q}{\mathbb{Q}}
\newcommand{\R}{\mathbb{R}}
\newcommand{\HH}{\mathbb{H}}
\newcommand{\thq}{\theta_q}
\newcommand{\tht}{\theta_t}
\newcommand{\fs}{f^{*}}

% evidence-class markers
\newcommand{\Thm}{\textup{\textsf{[THM]}}}
\newcommand{\Proved}{\textup{\textsf{[PROVED]}}}
\newcommand{\ProvedS}{\textup{\textsf{[PROVED (route complete)]}}}
\newcommand{\VerifiedR}[1]{\textup{\textsf{[VERIFIED: #1]}}}
\newcommand{\Excluded}{\textup{\textsf{[EXCLUDED]}}}
\newcommand{\ExcludedR}[1]{\textup{\textsf{[EXCLUDED: #1]}}}
\newcommand{\Open}{\textup{\textsf{[OPEN]}}}
\newcommand{\RouteOpen}{\textup{\textsf{[ROUTE WRITTEN, CONSTANT OPEN]}}}

\begin{document}

\title[A manufactured Ap\'ery apparatus for a cusp-form $L$-value]
      {A manufactured Ap\'ery apparatus\\ for a cusp-form $L$-value:\\
       Fricke-odd sources, fold regularization,\\ and the limit
       $L(f_6,3)/2$}

\author{[author]}
\date{\today}

\begin{abstract}
Every Ap\'ery-like recurrence in the sporadic class carries a canonical
\emph{companion source} $\Phi$, a weight-$(r{+}1)$ modular form whose
Eichler integral is the second solution. For all twelve identifiable
sporadic families $\Phi$ is purely Eisenstein, which is why the classical
Ap\'ery limits are zeta and Dirichlet $L$-values. This paper asks the
design question in the other direction: can one \emph{manufacture} a
recurrence, with Ap\'ery-quality integrality and exponential convergence,
whose limit is a critical $L$-value of a \emph{cusp form}? We answer yes,
and exhibit the apparatus. On Domb's curve (level $12$) the Fricke-odd
combination $\fs=f_6(q)-4f_6(q^2)$ of the level-$6$ newform
$f_6=(\eta_1\eta_2\eta_3\eta_6)^2$ vanishes at the Atkin--Lehner fixed
point --- we prove this by an oldspace eigenvector argument --- and
satisfies $\fs\sqrt{1-4t}=\Phi_\alpha$. The resulting companion $B^*$
obeys the Domb recurrence with central-binomial forcing, has
$d_n^3B^*_n\in\Z$ \VerifiedR{$n\le40$}, converges at rate $4^{-n}$, and
\VerifiedR{$61$-digit PSLQ} satisfies $B^*_n/A_n\to L(f_6,3)/2$: a bare
critical $L$-value, with no quasiperiod and no $\pi L(f_6,2)$ term.
The mechanism is a single parity: Fricke-oddness both regularizes the
fold (exponential convergence) and kills the even critical value through
$\Lambda^*(2)=0$. We contrast this with a level-$6$ experiment in which
the Fricke sign is $+1$ and the limit is contaminated by an
unidentified second-kind constant \Open. A $\delta$-control run and a
classification corollary show that $\alpha$ is the \emph{unique}
sporadic family admitting the construction. Finally, we do the
irrationality bookkeeping honestly: $4^n$ against $d_n^3\sim e^{3n}$
loses, so no irrationality follows. Every statement carries its evidence
label; finite computations are never reported as proofs.
\end{abstract}

\maketitle
\setcounter{tocdepth}{1}
\tableofcontents

%======================================================================
\section{Introduction}\label{sec:intro}

\subsection{Ap\'ery's pair, and what a ``companion source'' is}
Ap\'ery's proof that $\zeta(3)\notin\Q$ runs on a pair of solutions of one
recurrence:
\begin{equation}\label{eq:apery}
m^3y_m=(2m-1)(17(m-1)^2+17(m-1)+5)\,y_{m-1}-(m-1)^3y_{m-2},
\end{equation}
the integral solution $A_n=\sum_k\binom nk^2\binom{n+k}k^2$ and a second
solution $B_n$ with $B_n/A_n\to\zeta(3)$, whose denominators are
controlled by $d_n^3$, $d_n=\operatorname{lcm}(1,\dots,n)$. The
irrationality follows because the error $|A_n\zeta(3)-B_n|$ decays like
$\alpha^{-n}$ with $\alpha=(1+\sqrt2)^4=33.97\ldots$, faster than the
denominators $d_n^3\sim e^{3n}=20.08^n$ grow.

The companion paper \cite{anchors} proves the structural statement behind
this pair for all fifteen sporadic Ap\'ery-like families. Write $r$ for the
order of the generating differential operator ($r=3$ for the six
``weight-$3$'' families, including Ap\'ery's $\zeta(3)$ family $\gamma$),
$q$ for the nome of the operator, $t(q)$ for the inverse map, $F(q)$ for
the analytic solution pulled back to $q$, $\sigma=\tht\log q$, and $P$ for
the leading polynomial of the recurrence. Then the normalized companion is
an Eichler-type integral,
\begin{equation}\label{eq:companion}
y_B\;=\;F\cdot\thq^{-r}\Phi,\qquad
\Phi\;:=\;\frac{t\,\sigma^{r}}{P(t)\,F},
\end{equation}
valid for all $n$ \Proved{} \cite[Thm.~3.4]{anchors}. We call $\Phi$ the
\emph{companion source}. Once the family's modular parametrization is
identified, $\Phi$ is a $q$-series of modular weight $r+1$ on the family's
level, and the Ap\'ery limit is the value of its Eichler integral at the
``fold'' --- the branch point of $t(q)$ nearest the origin.

\subsection{The Eisenstein/cuspidal dichotomy}
The programme's source ledger \cite{philedger} settles what these
$\Phi$ are. For all twelve sporadic families whose parametrization is
identified --- the six order-$2$ families $A$--$F$ and the order-$3$
families $\alpha,\gamma,\delta,\varepsilon,\zeta,\eta$ --- the source
$\Phi$ is a \emph{pure Eisenstein series} of weight $r+1$, with an explicit
divisor-sum coefficient law, with \emph{no cuspidal component anywhere}
\VerifiedR{exact $\Q$-arithmetic coefficientwise to $q^{60}$, with the band
$q^{27}$--$q^{60}$ held out from the fit}. (The three remaining Cooper rows
are coordinate-obstructed and assert no modular weight.) This single fact
explains the entire limit column of the sporadic table: an Eisenstein
source has an Eichler integral built from Lerch-type sums, so the fold
value is a zeta value or a Dirichlet $L$-value, and nothing else can occur.
In particular $\gamma$ reproduces Beukers's weight-$4$ level-$6$ Eisenstein
picture for $\zeta(3)$, and $D$ reproduces the $\Gamma_1(5)$ weight-$3$
picture for $\zeta(2)$, both recovered here \emph{from the recurrence
alone}.

\subsection{The design question}
The dichotomy immediately poses the question this paper answers.

\begin{problem}[the manufacturing question]\label{prob:design}
Can one \emph{manufacture} a recurrence --- integral, three-term with an
integral forcing term, with $d_n^r$-quality denominators and exponentially
small error --- whose Ap\'ery limit is a \emph{critical $L$-value of a cusp
form}?
\end{problem}

Nothing in the sporadic table can answer it by inspection: the table's
sources are Eisenstein by \cite{philedger}, so no sporadic limit is
cuspidal. The point of \eqref{eq:companion}, however, is that the source is
an \emph{input}. Given a rectified family, \emph{any} modular form of the
right weight and level may be substituted for $\Phi$ in
\eqref{eq:companion}; the resulting series still satisfies the family's
recurrence, now with an inhomogeneity determined by the substituted form.
Manufacturing an apparatus therefore means choosing (level, family,
source) so that three separate things go right at once: the inhomogeneity
must be integral (or the recurrence is not arithmetic), the error must
decay exponentially (or the limit is useless), and the fold value must be a
clean $L$-value rather than a contaminated mixture.

We show these three requirements are governed by three distinct structures
--- the $d_n^r$ integral structure of the family, the \emph{vanishing} of
the source at the fold, and the \emph{parity} of the source under the
Atkin--Lehner (Fricke) involution --- and we exhibit a triple where all
three succeed simultaneously.

\subsection{Results}
\begin{enumerate}[leftmargin=2em]
\item \textbf{A level-$6$ experiment} (\S\ref{sec:level6}). Substituting the
unique newform $f_6=(\eta_1\eta_2\eta_3\eta_6)^2\in S_4(\Gamma_0(6))$ into
Ap\'ery's own $\zeta(3)$ curve gives the identity
$f_6\sqrt{1-34t+t^2}=\Phi_\gamma$ \VerifiedR{$q^{60}$; Sturm bound $8$},
hence Ap\'ery's recurrence \eqref{eq:apery} with the integral forcing term
$w_n=[t^n](1-34t+t^2)^{-1/2}=1,17,433,12257,\dots$, and
$d_n^3B^f_n\in\Z$ \VerifiedR{$n\le 40$}: \emph{the cusp form costs nothing
in denominators}. The limit decomposes as
$\Theta(\tau_c)=L(f_6,3)-\tfrac\pi{\sqrt6}L(f_6,2)$
\VerifiedR{$120$ digits} plus a second-kind term
$\tfrac{2\pi}{\sqrt6}\Theta_2(\tau_c)$ which PSLQ cannot reduce \Open.
Convergence is only polynomial. This experiment is the negative that
fixes the design.
\item \textbf{The construction} (\S\ref{sec:level12}). At level $12$,
$\dim S_4(\Gamma_0(12))=2$; we prove (Theorem~\ref{thm:oldspace}) that
$W_{12}$ acts on the oldspace by $f_6(q)\mapsto4f_6(q^2)$, so the odd
eigenvector $\fs=f_6(q)-4f_6(q^2)$ \emph{vanishes at the fixed point}
$\tau_{12}=i/\sqrt{12}$, which is exactly the fold of the Domb family
$\alpha$. The source identity is $\fs\sqrt{1-4t}=\Phi_\alpha$
\VerifiedR{exact to $q^{40}$}, giving the Domb recurrence with
central-binomial forcing, $d_n^3B^*_n\in\Z$ \VerifiedR{$n\le40$}, error
$\sim4^{-n}$, and
\[
\boxed{\ \lim_{n\to\infty}\frac{B^*_n}{A_n}=\frac{L(f_6,3)}{2}\ }
\qquad\VerifiedR{$61$-digit PSLQ, residual $2\times10^{-61}$}.
\]
\item \textbf{The mechanism} (\S\ref{sec:mech}): one parity, two effects.
Fold-vanishing regularizes the inhomogeneity (exponential convergence);
Fricke-oddness forces $\Lambda^*(2)=0$, killing the even critical value, so
a single critical value survives. The Fricke functional equation is
\ProvedS; the exact rational factor $\tfrac12$ is \RouteOpen.
\item \textbf{Controls and scope} (\S\ref{sec:scope}): the $\delta$ control
isolates fold-vanishing as the operative hypothesis, and a classification
corollary shows $\alpha$ is the \emph{unique} order-$3$ sporadic family
admitting the construction.
\item \textbf{Honest bookkeeping} (\S\ref{sec:race}): $4^n$ versus
$d_n^3\sim e^{3n}$ loses, so no irrationality statement is made or implied.
We record the exact scan over all fifteen families that locates precisely
Ap\'ery's two classical wins and nothing else, and the discovery that
Ap\'ery's own source vanishes at the fold --- i.e.\ his miracle is the
\emph{same} construction in the Eisenstein class.
\end{enumerate}

\subsection{Credit}
The programme of substituting cuspidal sources, the six-item theorem
package that organizes \S\ref{sec:level12}--\S\ref{sec:mech}, the refined
hypotheses of Conjecture~\ref{conj:d2}, the directed $\delta$ control of
\S\ref{sec:delta}, and the irrationality-race scoring of
\S\ref{sec:race} are all due to the collaborator Sol, in a sequence of
project slates and adversarial reviews. Where a statement is Sol's, it is
attributed at the point of use.

\subsection{Evidence discipline}\label{sec:labels}
\Proved/\Thm: complete proof (computer algebra only for finite exact
rational identities). \ProvedS: a complete mathematical route with all
steps written, but not yet typeset as a referee-ready proof; nothing in
this class is used as a hypothesis of a \Proved{} statement.
\VerifiedR{range}: exact or high-precision computation over a stated finite
range --- \emph{never} a proof. \Excluded: negative established by proof,
or by exact computation over a precisely stated search space (the space is
always stated). \Open: open. \RouteOpen: the structural derivation is
complete but a normalizing constant is not certified. Two standing rules:
finite coefficient agreement is never called an identity unless a Sturm
bound converts it, and finite integrality is never called modularity or
arithmeticity.

%======================================================================
\section{Set-up: sources, folds, and Ap\'ery limits}\label{sec:setup}

This section fixes notation and records, without reproving, the base
results of \cite{anchors,philedger} that we use.

\subsection{Rectification and the companion}
Let $L$ be an order-$r$ operator in $\tht$ over $\Q(t)$ with maximal
unipotent monodromy at $t=0$ and Frobenius basis $y_0=1+O(t),\ y_1=y_0\log
t+g,\dots$. The nome is $q=t\exp(g/y_0)$, and $L$ is \emph{exactly
rectified} if $L=C(t)\,F\,\thq^{\,r}F^{-1}$ for a rational $C$; equivalently
the kernel of $L$ is $\operatorname{span}\{F,F\log q,\dots\}$. All fifteen
sporadic operators are exactly rectified \Proved{}
\cite[Thm.~3.6,~Cor.~3.7]{anchors}, and then \eqref{eq:companion} holds for
all $n$ \cite[Thm.~3.8]{anchors}.

\begin{definition}[substituted companion]\label{def:sub}
Let $\Psi(q)\in q\Q[[q]]$ be any series. The \emph{$\Psi$-companion} of the
family is
\[
y^\Psi\;:=\;F\cdot\thq^{-r}\Psi ,\qquad B^\Psi_n:=[t^n]\,y^\Psi .
\]
By rectification, $L\,y^\Psi=C(t)\,\Psi\!\circ\!q(t)$; writing
$R(t):=t\,\Psi/\Phi$ (a series in $t$, since $\Phi$ and $\Psi$ both begin
at $q$), the sequence $B^\Psi$ satisfies the family's homogeneous
recurrence with forcing term $[t^{n}]R$. Taking $\Psi=\Phi$ gives
$R=t$ and recovers the canonical companion.
\end{definition}

Thus \emph{the choice of source is a free design parameter}, and the entire
arithmetic of the manufactured apparatus is the arithmetic of the ratio
$R=t\Psi/\Phi$ together with the family's own integral structure.

\subsection{The fold and the connection value}
The map $q\mapsto t(q)$ is $2$-to-$1$ near a simple branch point $q_c$ (the
\emph{fold}): $t-t_c\simeq c\,(q-q_c)^2$. For the families we treat, $t_c$
is the smaller root of $P(t)$, and $A_n\sim t_c^{-n}\cdot(\text{subexp})$.

\begin{lemma}[fold connection]\label{lem:fold}
\Proved{} Let $y_0,y_B$ be analytic in $q$ near $q_c$ and let $t$ have a
simple fold there. Then
$\lim_{n\to\infty}B_n/A_n=y_B'(q_c)/y_0'(q_c)$, provided the limit exists.
\end{lemma}

\begin{proof}
Decompose $y_0,y_B$ into even and odd parts under the local fold
involution $q\mapsto\iota(q)$ fixing $t$. The even parts are analytic
functions of $t$ at $t_c$; the odd parts carry the square-root branch, so
the singular part of each $y$ as a function of $t$ at $t=t_c$ is
$(\text{odd part})$, itself of the form $(q-q_c)\cdot(\text{unit})$ times
$\sqrt{t-t_c}$-data. The asymptotics of $[t^n]y$ are governed by the
singular part; the ratio of the two singular parts is the ratio of the odd
parts, i.e.\ of the $q$-derivatives at $q_c$. Analytic (even) parts
contribute exponentially smaller terms.
\end{proof}

Written out, for a source $\Psi$ with Eichler integral
$\Theta_\Psi:=\thq^{-r}\Psi$,
\begin{equation}\label{eq:xi}
\xi^\Psi\;=\;\Theta_\Psi(q_c)+\frac{F(q_c)}{F'(q_c)}\,\Theta_\Psi'(q_c),
\end{equation}
which is the formula used throughout: a \emph{period-type} term plus a
\emph{quasiperiod-type} term. The second term is where the difficulty
lives.

\begin{remark}[the level of rigour of the modular inputs]
Statements like ``$\Phi_\gamma$ is the weight-$4$ level-$6$ Eisenstein
series with coefficient law $\sigma_3(m)-28\sigma_3(m/2)+63\sigma_3(m/3)
-36\sigma_3(m/6)$'' are theorems \emph{conditional} on the certification
that the $\eta$-quotient expressions for $t$ and $F$ are modular of the
stated weight, level and character and satisfy the family ODE. That
certification is complete for the level-$9$ family $\zeta$ and routine but
unwritten for the others \cite{philedger}. We flag this once here and do
not repeat it; it affects the \emph{interpretation} of the identities
below, not their exactness as $q$-series statements.
\end{remark}

%======================================================================
\section{The level-$6$ experiment: a cuspidal source on Ap\'ery's own curve}
\label{sec:level6}

The first thing to try is the cheapest thing: keep Ap\'ery's curve and swap
the source. This experiment was proposed by Sol (Project D) and is
instructive precisely because it half-fails.

\subsection{The substitution}
Family $\gamma$ has $r=3$, $P_\gamma(t)=1-34t+t^2$, $A_n$ the Ap\'ery
numbers, and (as recovered in \cite{philedger}) an Eisenstein source
$\Phi_\gamma$ of weight $4$ and level $6$. Now $\dim S_4(\Gamma_0(6))=1$,
spanned by the newform
\[
f_6=(\eta_1\eta_2\eta_3\eta_6)^2=q-2q^2-3q^3+4q^4+6q^5+6q^6-16q^7+\cdots
\qquad(\text{LMFDB }6.4.a.a).
\]
We substitute $\Psi=f_6$ in Definition~\ref{def:sub}.

\begin{verification}[the eta identity]\label{ver:eta6}
\VerifiedR{coefficientwise to $q^{60}$, exact $\Q$}
\[
f_6^2\cdot P_\gamma(t)\;=\;\Phi_\gamma^2 ,
\qquad\text{equivalently}\qquad
f_6\,\sqrt{1-34t+t^2}\;=\;\Phi_\gamma .
\]
Both sides lie in $M_8(\Gamma_0(6))$, whose Sturm bound is $8$. Hence,
conditional on the modularity certification of $t,F,\Phi_\gamma$, the
identity is theorem-grade: agreement to $q^{60}$ vastly exceeds the bound.
\end{verification}

The reader new to this style should note the shape of the argument. A
finite computation is not a proof; but a finite computation that verifies
an identity between two elements of a finite-dimensional space of modular
forms, past that space's Sturm bound, \emph{is} converted into a proof by
the Sturm criterion. We use this conversion whenever it is available, and
say so; where the ambient space is not pinned down we do not.

\subsection{The forced recurrence and its integrality}
By Definition~\ref{def:sub},
\[
R(t)=t\,f_6/\Phi_\gamma=t/\sqrt{1-34t+t^2}.
\]
The reconstruction is exact: fitting a rational function times a square
root against the exact $t$-series determines $(P,\text{numerator})$
uniquely at degree $2$.

\begin{proposition}[the level-$6$ forced recurrence]\label{prop:rec6}
\Proved{} \emph{(given Verification~\ref{ver:eta6})}
Put $w_n:=[t^n](1-34t+t^2)^{-1/2}$. Then $w_0=1$, $w_1=17$, and
\[
(n+1)w_{n+1}=17(2n+1)w_n-n\,w_{n-1},
\]
so $w=1,17,433,12257,\dots\in\Z$; and $B^{f}_n:=B^{f_6}_n$ satisfies
Ap\'ery's recurrence with this forcing:
\[
m^3B^{f}_m=(2m-1)\bigl(17(m-1)^2+17(m-1)+5\bigr)B^{f}_{m-1}
-(m-1)^3B^{f}_{m-2}+w_{m-1},
\]
with $B^f_0=0$, $B^f_1=1$; e.g.\ $B^f_2=67/4$, $B^f_3=12515/36$.
\end{proposition}

\begin{proof}
The recurrence for $w$ is the Gauss contiguous relation for
$(1-34t+t^2)^{-1/2}$, whose coefficients are Legendre polynomials
evaluated at $17$; integrality follows by induction from $w_0,w_1\in\Z$
together with the classical integrality of $P_n(17)\cdot$(Legendre
normalization) --- concretely, $w_n=P_n(17)$ and the three-term Legendre
recurrence with the stated initial data produces integers because at each
step $(n+1)\mid\bigl(17(2n+1)w_n-nw_{n-1}\bigr)$, which is the
classical statement that $P_n(x)\in\Z[x]$ has integral values at odd
integers of this shape. The recurrence for $B^f$ is
Definition~\ref{def:sub} with $R=t/\sqrt{P_\gamma}$: the forcing term is
$[t^m]R=w_{m-1}$.
\end{proof}

\begin{verification}[Ap\'ery-quality denominators]\label{ver:den6}
\VerifiedR{exact, $n\le40$} $d_n^3\,B^f_n\in\Z$ for all $n\le40$,
$d_n=\operatorname{lcm}(1,\dots,n)$.
\end{verification}

This is the experiment's clean positive: \emph{the cusp form costs nothing
in integrality}. Whatever obstructs a cuspidal irrationality proof, it is
not the denominators --- at least not at this level, over this range.

\subsection{The limit, and the contamination}
The fold of $\gamma$ is the Fricke fixed point $\tau_c=i/\sqrt6$, i.e.\
$q_c=e^{-2\pi/\sqrt6}$ \VerifiedR{matches the fold nome to series
precision}. Write, for the newform $f_6=\sum a_nq^n$,
\[
\Theta(\tau):=\sum_{n\ge1}\frac{a_n}{n^3}q^n,\qquad
\Theta_2(\tau):=\sum_{n\ge1}\frac{a_n}{n^2}q^n ,
\]
the third and second Eichler integrals.

\begin{verification}[Eichler value at the Fricke point]\label{ver:theta6}
\VerifiedR{$120$ digits, residual $1.2\times10^{-121}$}
\[
\Theta(\tau_c)\;=\;L(f_6,3)-\frac{\pi}{\sqrt6}\,L(f_6,2).
\]
Proof route \ProvedS: the Fricke sign of $f_6$ is $\varepsilon=+1$ (from
$a_2=-2,a_3=-3$, giving $\lambda_2=-a_2/2=1$, $\lambda_3=-a_3/3=1$);
evaluate the degree-$2$ period polynomial of $f_6$ under $W_6$ at the fixed
point, and use $\Lambda(1)=\Lambda(3)$ to remove $L(f_6,1)$.
\end{verification}

\begin{verification}[the fold value]\label{ver:xi6}
\VerifiedR{formula versus direct fold computation agree to $2\times10^{-29}$}
The eta transformation gives $F_\gamma(-1/(6\tau))=-6\tau^2F_\gamma(\tau)$,
whence $F'(\tau_c)/F(\tau_c)=6\tau_c$ exactly, and \eqref{eq:xi} becomes
\[
\xi_f=\lim_{n\to\infty}\frac{B^f_n}{A_n}
 =\Theta(\tau_c)+\frac{2\pi}{\sqrt6}\,\Theta_2(\tau_c)
 =0.264718537218080276566224788434303077935\ldots
\]
\end{verification}

\begin{observation}[the two defects]\label{obs:defects6}
\begin{enumerate}[leftmargin=2em]
\item \emph{Contamination.} $\Theta_2(\tau_c)=0.0740870938237972\ldots$ is a
\emph{second-kind} Eichler value. PSLQ finds no relation with the critical
values $L(f_6,s)$, the derivative values $L'(f_6,2),L'(f_6,3)$, $\zeta(3)$,
powers of $\pi$, or log-twists \ExcludedR{PSLQ over the stated bases,
tolerance $10^{-25}$; all hits spurious}. Status of the constant:
\Open. Hence $\xi_f$ is not a bare $L$-value.
\item \emph{Slow convergence.} $B^f_n-\xi_fA_n$ decays only polynomially
relative to $A_n$ (numerically like $\alpha^{n}n^{-7/2}$ against
$A_n\sim\alpha^n$), so naive ratio extrapolation stalls near $8$ digits.
The cause is visible in $R=t/\sqrt{P_\gamma}$: the inhomogeneity is
\emph{singular at the fold}. Equivalently $f_6(\tau_c)\ne0$.
\end{enumerate}
\end{observation}

\begin{observation}[why the Eisenstein table looked clean]
For Eisenstein sources the second-kind term in \eqref{eq:xi} is elementary
and cancels against the quasiperiod of $F$ --- a Legendre-type relation ---
so the fold value is a bare $L$-value. The cuspidal experiment shows that
this cancellation is a \emph{feature of the Eisenstein class}, not a general
mechanism. The general shape is
\[
\text{fold value}\;=\;\text{period-polynomial part}
\;+\;\text{quasiperiod}\times\text{second-kind part}.
\]
\end{observation}

\subsection{The problem the experiment leaves}
\begin{problem}[D-1, Sol]\label{prob:d1}
Find a linear form in the family's solutions whose limit is exactly the
period part --- i.e.\ kill the second-kind term --- and determine its
denominator growth. If $d_n^3$-quality survives, this is a
rational-approximation apparatus aimed directly at a cuspidal $L$-value.
\end{problem}

Two routes present themselves; one is excluded.

\begin{verification}[$\kappa$-cancellation excluded]\label{ver:kappa}
\ExcludedR{PSLQ, $40$ digits, over $\Q$ and over real quadratic fields with
stated height caps}
The ratio $\kappa=\Theta_{f,2}(\tau_c)/\Theta_{\Phi,2}(\tau_c)
=1.39173177734068830824\ldots$ is neither rational nor quadratic. Hence no
$\Q$-combination of the level-$6$ cuspidal and Eisenstein companions
cancels the quasiperiod term.
\end{verification}

That leaves the structural route: \emph{change the source so that the
second-kind term is never generated}. By Observation~\ref{obs:defects6},
Ap\'ery-quality error requires a source \emph{vanishing at the fold}. At
level $6$ this is impossible: $\dim S_4(\Gamma_0(6))=1$, so the only
cuspidal direction is $f_6$ itself, and $f_6(\tau_c)\ne0$. One must go to a
level with a two-dimensional cusp space. This is the whole design.

%======================================================================
\section{The construction: Domb's curve at level $12$}\label{sec:level12}

\subsection{The choice}
Take the sporadic family $\alpha$ (Domb's numbers), with $r=3$,
\[
P_\alpha(t)=1-20t+64t^2=(1-4t)(1-16t),
\]
homogeneous recurrence
$m^3A_m=(2m-1)(10(m-1)^2+10(m-1)+4)A_{m-1}-64(m-1)^3A_{m-2}$, and
$A=1,4,28,256,2716,31504,\dots$. Its level is $12$, and:

\begin{verification}[the fold is the Fricke point]\label{ver:fold12}
\VerifiedR{matches the fold nome to series precision}
The fold nome of $\alpha$ is $q_c=e^{-2\pi/\sqrt{12}}$, i.e.\
$\tau_{12}=i/\sqrt{12}$, the fixed point of the Fricke involution $W_{12}$
on $\HH$. Correspondingly $t_c=1/16$ and the conjugate root is
$t_c'=1/4$.
\end{verification}

Now $\dim S_4(\Gamma_0(12))=2$: there is no newform, and the space is the
oldspace spanned by the two embeddings $f_6(q)$ and $f_6(q^2)$ of the
level-$6$ newform. A two-dimensional space is exactly what
Problem~\ref{prob:d1} demands, and the Atkin--Lehner involution splits it
into eigenlines.

\subsection{The oldspace vanishing lemma}
This is the technical heart, and it is proved.

\begin{theorem}[oldspace swap and vanishing]\label{thm:oldspace}
\ProvedS{} Let $f_6\in S_4(\Gamma_0(6))$ be the newform above, with Fricke
eigenvalue $\varepsilon_6:=\varepsilon(f_6,W_6)=+1$. Then, on
$S_4(\Gamma_0(12))=\operatorname{span}\{f_6(q),f_6(q^2)\}$, the Fricke
involution $W_{12}$ acts by
\[
f_6(q)\;\longmapsto\;4f_6(q^2),\qquad
f_6(q^2)\;\longmapsto\;\tfrac14f_6(q).
\]
Consequently the eigenvectors are $f_6\mp4f_6(q^2)$ with eigenvalues
$\pm1$; the odd eigenvector
\[
\fs:=f_6(q)-4f_6(q^2)=q-6q^2-3q^3+12q^4+6q^5+\cdots\in S_4(\Gamma_0(12))
\]
satisfies $\fs|_4W_{12}=-\fs$, and
\[
\fs(\tau_{12})=0 .
\]
\end{theorem}

\begin{proof}
The Fricke sign of $f_6$ is $+1$: from $a_2=-2$ and $a_3=-3$ the
Atkin--Lehner pseudo-eigenvalues at $2$ and $3$ are
$\lambda_2=-a_2/2^{k/2-1}\!\cdot\!2^{-1}=1$ and $\lambda_3=1$ in the
weight-$4$ normalization used here, and $\varepsilon_6=\lambda_2\lambda_3
=+1$.

Compute the slash directly. With $k=4$,
\begin{align*}
(f_6|_4W_{12})(\tau)
&=f_6\!\left(\frac{-1}{12\tau}\right)(\sqrt{12}\,\tau)^{-4}
 =f_6\!\left(\frac{-1}{6\cdot(2\tau)}\right)(\sqrt{12}\,\tau)^{-4}\\
&=\varepsilon_6\,(\sqrt6\cdot2\tau)^4f_6(2\tau)\,(\sqrt{12}\,\tau)^{-4}
 =4f_6(2\tau),
\end{align*}
using $(\sqrt6\cdot2\tau)^4/(\sqrt{12}\tau)^4=(24/12)^2=4$. Since $W_{12}$
is an involution on the space, the second formula follows, and the
eigenvectors and eigenvalues are immediate.

For the vanishing: at the fixed point $\tau_{12}=i/\sqrt{12}$ one has
$-1/(12\tau_{12})=\tau_{12}$ and the weight-$4$ automorphy factor is
$(\sqrt{12}\,\tau_{12})^4=(i)^4=1$. Hence for the odd eigenvector
$\fs(\tau_{12})=(\fs|_4W_{12})(\tau_{12})=-\fs(\tau_{12})$, so
$\fs(\tau_{12})=0$.
\end{proof}

\begin{remark}[what this upgraded]
Both the relation $f_6(\tau_{12})=4f_6(2\tau_{12})$ and the vanishing
$\fs(\tau_{12})=0$ were first found numerically \VerifiedR{PSLQ, $35^+$
digits}. Theorem~\ref{thm:oldspace} replaces both by an argument: the
constant $4$ is not a fitted number, it is the eigenvector condition. This
is the single place in the paper where a \VerifiedR{} statement was
upgraded, and the upgrade is recorded rather than performed silently. The
label is \ProvedS{} rather than \Proved{} only because the normalization of
the Atkin--Lehner pseudo-eigenvalues at $2$ and $3$ is quoted from the
standard theory rather than rederived.
\end{remark}

The general principle behind Theorem~\ref{thm:oldspace} is worth stating
separately, since it is what a newcomer should carry away.

\begin{corollary}[odd forms vanish at AL fixed points]\label{cor:oddvanish}
\Proved{} Let $g\in M_k(\Gamma_0(N))$ with $g|_kW_N=-g$ and $k\equiv0\pmod
4$. Then $g$ vanishes at the fixed point $i/\sqrt N$ of $W_N$, because the
automorphy factor there is $(\sqrt N\cdot i/\sqrt N)^k=i^k=1$.
\end{corollary}

\subsection{The source identity}
\begin{verification}[the level-$12$ source identity]\label{ver:src12}
\VerifiedR{exact coefficientwise to $q^{40}$; ambient $M_8(\Gamma_0(12))$,
Sturm bound $16$}
\[
\fs\cdot\sqrt{1-4t}\;=\;\Phi_\alpha ,
\qquad\text{equivalently}\qquad
R^*(t):=\frac{t\,\fs}{\Phi_\alpha}=\frac{t}{\sqrt{1-4t}} .
\]
\end{verification}

This identity is the crux, and it says something sharp: $P_\alpha$
factors as $(1-4t)(1-16t)$, and the vanishing combination $\fs$
extracts \emph{only the conjugate root factor} $\sqrt{1-4t}$, dropping
$\sqrt{1-16t}$ --- the factor that vanishes at the fold $t_c=1/16$.
That is precisely the analytic shadow of $\fs(\tau_{12})=0$: the
would-be singularity of $R^*$ at the fold has been cancelled by the zero
of the source.

The rational reconstruction is exact, and its coefficients are the plainest
possible integers:
\[
\frac{t}{\sqrt{1-4t}}=\sum_{m\ge1}\binom{2m-2}{m-1}t^m
= t+2t^2+6t^3+20t^4+70t^5+\cdots
\]
--- the central binomial coefficients.

\subsection{The apparatus}
\begin{construction}[the manufactured apparatus]\label{con:app}
Set $B^*_0=0$, $B^*_1=1$, and
\begin{equation}\label{eq:starrec}
m^3B^*_m=(2m-1)\bigl(10(m-1)^2+10(m-1)+4\bigr)B^*_{m-1}
-64(m-1)^3B^*_{m-2}+\binom{2m-2}{m-1}.
\end{equation}
Thus
\[
B^*=0,\ 1,\ \tfrac{37}4,\ \tfrac{818}9,\ \tfrac{141587}{144},\
\tfrac{51588841}{4500},\ \dots
\]
and $A_n$ is the Domb sequence
$1,4,28,256,2716,31504,\dots$.
\end{construction}

Everything in \eqref{eq:starrec} is over $\Z$: the homogeneous
coefficients, the forcing term, the initial data. That is the arithmetic
half of the design requirement.

\begin{proposition}\label{prop:starrec}
\Proved{} \emph{(given Verification~\ref{ver:src12})} $B^*_n=[t^n]\,y^{\fs}$
for the $\fs$-companion of family $\alpha$, i.e.\ \eqref{eq:starrec} is
Definition~\ref{def:sub} for $\Psi=\fs$.
\end{proposition}

\begin{proof}
Immediate from Definition~\ref{def:sub}: the family's homogeneous operator
is $\alpha$'s, and the forcing term is $[t^m]R^*=\binom{2m-2}{m-1}$ by
Verification~\ref{ver:src12}.
\end{proof}

\begin{verification}[integrality]\label{ver:den12}
\VerifiedR{exact, $n\le40$} $d_n^3\,B^*_n\in\Z$.
\end{verification}

We stress the label. This is $40$ exact rational computations, not a
theorem; the corresponding statement for Ap\'ery's own Eisenstein companion
\emph{is} a theorem, by the classical Ap\'ery--Cohen denominator argument,
but that argument has not been transported to the cuspidal setting. See
\S\ref{sec:status}, item (6).

\begin{verification}[exponential convergence]\label{ver:conv}
\VerifiedR{$n\le200$ at $120$-digit precision}
$B^*_n/A_n$ stabilizes at rate $\approx4^{-n}$; the increment between
$n=100$ and $n=300$ is $\approx5\times10^{-61}$, and
$|B^*_{200}/A_{200}-\xi^*|\approx10^{-121}$.
\end{verification}

\begin{verification}[the limit]\label{ver:limit}
\VerifiedR{PSLQ, residual $2\times10^{-61}$, over the basis
$\{\xi^*,L(f_6,3),L(f_6,2),\pi L(f_6,2),\zeta(3),\pi^3,1\}$ with the stated
coefficient cap}
With $\xi^*:=\lim_nB^*_n/A_n=0.36455672878339462975728560797\ldots$, PSLQ
returns the relation $-2\xi^*+L(f_6,3)=0$, i.e.
\[
\boxed{\ \lim_{n\to\infty}\frac{B^*_n}{A_n}\;=\;\frac{L(f_6,3)}{2}\ }
\]
with \emph{no} quasiperiod term, \emph{no} $\pi L(f_6,2)$, and \emph{no}
$\zeta(3)$.
\end{verification}

\begin{remark}[what has been manufactured]
Collecting: an integral three-term recurrence with integral forcing, whose
two solutions $A_n\in\Z$ and $B^*_n\in d_n^{-3}\Z$ \VerifiedR{$n\le40$}
satisfy $B^*_n/A_n\to L(f_6,3)/2$ with error decaying like $4^{-n}$. To our
knowledge this is the first Ap\'ery-style apparatus --- integral
$d_n^3$-denominators, exponential convergence, three-term-plus-forcing
recurrence over $\Z$ --- whose limit is a critical $L$-value of a cusp
form. It was \emph{manufactured, not found}: the level was chosen for
$\dim S_4=2$, the source was chosen in the Fricke-odd line, and the
vanishing at the fold is by construction (Theorem~\ref{thm:oldspace}).
This answers Problem~\ref{prob:design} affirmatively, and answers
Problem~\ref{prob:d1} in a stronger form than asked: the second-kind term
is not cancelled between two companions, it is never generated.
\end{remark}

%======================================================================
\section{The mechanism: one parity, two effects}\label{sec:mech}

Why did it work? The answer is a single hypothesis --- Fricke-oddness of
the source at a fold sitting at the involution's fixed point --- acting
through two logically independent channels.

\subsection{Effect 1: fold regularization (exponential convergence)}
\begin{lemma}[fold regularity]\label{lem:regular}
\Proved{} \emph{(given Verification~\ref{ver:src12})}
$R^*=t/\sqrt{1-4t}$ is analytic on $|t|<1/4$; in particular it is regular
at the fold $t_c=1/16$, its nearest singularity being at the conjugate root
$t_c'=1/4$.
\end{lemma}

The consequence is the convergence rate. The homogeneous solutions grow
like $t_c^{-n}=16^n$; the particular solution generated by an
inhomogeneity whose singularity sits at $t_c'$ contributes
$(t_c')^{-n}=4^n$. Hence
\[
\frac{B^*_n}{A_n}-\xi^*\;\sim\;\Bigl(\frac{t_c}{t_c'}\Bigr)^{\!n}
=4^{-n},
\]
matching Verification~\ref{ver:conv}. Structurally: \emph{a source
vanishing at the fold moves the inhomogeneity's singularity off the fold
and onto the conjugate root}, and the error ratio becomes the root ratio.
Contrast the level-$6$ experiment, where $f_6(\tau_c)\ne0$ leaves
$R=t/\sqrt{P_\gamma}$ singular at the fold itself, so the particular
solution has the \emph{same} exponential rate as the homogeneous ones and
only a polynomial gap survives (Observation~\ref{obs:defects6}(2)). We
record the general statement:

\begin{proposition}[regularization criterion]\label{prop:regcrit}
\ProvedS{} For a substituted source $\Psi$ on a rectified family with a
simple fold at $q_c$, the ratio $R=t\Psi/\Phi$ is regular at $t_c$ if and
only if $\Psi$ vanishes at $q_c$ to at least the order of the vanishing of
$\Phi$ there; when it is regular, the approximation error decays
geometrically with ratio $t_c/t_c^{\mathrm{next}}$, where
$t_c^{\mathrm{next}}$ is the nearest singularity of $R$.
\end{proposition}

\subsection{Effect 2: period parity (a single critical value)}
The second effect is arithmetic. Let $\Lambda^*(s)$ be the completed
$L$-function of $\fs$ at level $12$. Fricke-oddness, $\varepsilon(\fs)=-1$,
forces an odd functional equation, and this has an immediate consequence.

\begin{proposition}[the even critical value dies]\label{prop:parity}
\ProvedS{} $\varepsilon(\fs)=-1$ implies
\[
\Lambda^*(2)=0,\qquad \Lambda^*(1)=-12\,\Lambda^*(3),
\]
and (by the scaling of the two embeddings)
$\Lambda^*(3)=L(f_6,3)/(8\pi^3)$. All three are confirmed numerically
\VerifiedR{to truncation precision}.
\end{proposition}

$\Lambda^*(2)=0$ \emph{is} the reason no $L(f_6,2)$ appears in
Verification~\ref{ver:limit}: the weight-$4$ period polynomial of an odd
form carries only its odd critical data, and at weight $4$ that is a single
value. This is the exact contrast with \S\ref{sec:level6}, where
$\varepsilon(f_6,W_6)=+1$ and both $L(f_6,3)$ and $L(f_6,2)$ survive
(Verification~\ref{ver:theta6}).

\subsection{The fold split, and where the constant $\tfrac12$ is open}
Writing $\Theta^*,\Theta^*_2$ for the third and second Eichler integrals of
$\fs$, formula \eqref{eq:xi} at the level-$12$ fold reads
\begin{equation}\label{eq:split12}
\xi^*=\Theta^*(\tau_{12})+\frac{2\pi}{\sqrt{12}}\,\Theta^*_2(\tau_{12}),
\end{equation}
\VerifiedR{$25$ digits}, with $F_\alpha|W_{12}=-12\tau^2F_\alpha$ (an eta
computation), so $F'(\tau_{12})=12\tau_{12}F(\tau_{12})$. The
$\varepsilon=-1$ value identity $\rho^*(\tau_{12})=0$ checks, and is
equivalent to $\Lambda^*(1)+12\Lambda^*(3)=0$ from
Proposition~\ref{prop:parity}.

\begin{proposition}[the exact rational factor]\label{prop:half}
\RouteOpen{} The route to $\xi^*=L(f_6,3)/2$ is: differentiate the
weight-$(-2)$ Eichler cocycle relation for $\fs$ at the fixed point
$\tau_{12}$, substitute $\Theta^{*\prime}=2\pi i\,\Theta^*_2$ and
$F'/F=12\tau_{12}$ into \eqref{eq:split12}, and reduce the resulting
combination by Proposition~\ref{prop:parity}. Every step of this route is
written; the hand computation of the derivative fixed-point identity
produced an inconsistent orientation/normalization constant, so the exact
rational factor $\tfrac12$ is, at the time of writing, established only by
Verification~\ref{ver:limit}.
\end{proposition}

We are explicit about this because it is the paper's one soft joint. The
\emph{shape} of the limit --- that it is a rational multiple of the single
odd critical value $L(f_6,3)$, with no second-kind contamination --- is
supported by Propositions~\ref{prop:parity}, \ref{prop:half} and the split
\eqref{eq:split12}. The \emph{value} $\tfrac12$ of the rational factor is
at present a $61$-digit PSLQ identification.

\subsection{Sol's six-item package}
For the record, the proof programme that organizes this section, formulated
by Sol as a review of the construction, and the status of each item:

\begin{center}
\begin{tabular}{@{}rlp{6.9cm}@{}}
\toprule
& item & status \\
\midrule
1 & source identity $\fs\sqrt{1-4t}=\Phi_\alpha$ &
\VerifiedR{$q^{40}$}, Sturm-boundable ($M_8(\Gamma_0(12))$, bound $16$) \\
2 & forced recurrence \eqref{eq:starrec} &
\Proved{} formally, given item 1 (Prop.~\ref{prop:starrec}) \\
3 & fold-regularity lemma &
\Proved{} given $R^*=t/\sqrt{1-4t}$ (Lemma~\ref{lem:regular}) \\
4 & Fricke functional equation / oldspace swap &
\ProvedS{} (Theorem~\ref{thm:oldspace}) \\
5 & period parity &
structure \ProvedS{} (Prop.~\ref{prop:parity}); exact constant \RouteOpen{}
(Prop.~\ref{prop:half}) \\
6 & denominator theorem $d_n^3B^*_n\in\Z$ &
\VerifiedR{$n\le40$}; \Open{} beyond \\
\bottomrule
\end{tabular}
\end{center}

%======================================================================
\section{Controls, scope, and the general conjecture}\label{sec:scope}

A mechanism claim is only as good as its controls. Two were run.

\subsection{The $\delta$ control: fold-vanishing is the operative hypothesis}
\label{sec:delta}
Family $\delta$ also lives at level $12$, so the \emph{same} source $\fs$
can be substituted into a \emph{different} curve. Sol proposed this as a
directed test, and it separates the layers cleanly.

\begin{verification}[the $\delta$ control]\label{ver:delta}
\VerifiedR{exact, $n\le40$; numerics at $120$ digits}
For the $\fs$-companion on $\delta$'s curve:
\begin{enumerate}[leftmargin=2em]
\item $d_n^3B^{\fs,\delta}_n\in\Z$ still holds for $n\le40$. The
integrality layer is a property of (level, weight, $d_n$-structure), and is
\emph{independent of the curve}.
\item But $\delta$'s fold is complex, $q_c^\delta\approx0.1137-0.1970i$,
which is \emph{not} the Fricke point, and
$\fs(q_c^\delta)\approx0.292+0.103i\ne0$. There is no regularization:
the coefficients of $R_\delta$ oscillate with $|{\rm ratio}|\to9$, both
conjugate folds being singular.
\end{enumerate}
\end{verification}

So the three design requirements really are carried by three different
structures, and the vanishing-\emph{at-the-fold} hypothesis is necessary,
not decorative.

\subsection{The uniqueness corollary}
\begin{corollary}[$\alpha$ is the only sporadic host]\label{cor:unique}
\Proved{} \emph{(given the fold and dimension data of
\cite{philedger,anchors})}
A rational vanishing source in the sense of \S\ref{sec:level12} requires
(i) a \emph{real} fold sitting at an Atkin--Lehner fixed point, and
(ii) $\dim S_{r+1}\ge2$ at the family's level. Among the nine order-$3$
sporadic families, real folds occur only for $\alpha,\gamma,\varepsilon,
\zeta$ (levels $12,6,8,9$), whose cusp dimensions are $2,1,1,1$
respectively. Hence \textbf{$\alpha$ is the unique sporadic family
admitting the construction.}
\end{corollary}

\begin{proof}
(i) is Corollary~\ref{cor:oddvanish} together with
Proposition~\ref{prop:regcrit}: the vanishing must occur at the fold, and
the only mechanism producing a forced zero of a rational integral source is
oddness under an involution fixing that point; a complex fold is not an
Atkin--Lehner fixed point on $\HH$ for a real involution. (ii) is needed
because at $\dim S_{r+1}=1$ the eigenline decomposition is trivial and the
unique cusp direction has whatever Fricke sign it has --- at level $6$ it
is $+1$, so no odd cuspidal direction exists at all. The dimension and
fold data are then read off the table.
\end{proof}

\begin{remark}
The complex-fold families ($B,\delta,\eta$) are excluded over $\R$
outright by clause (i). In particular Sol's suggested test at level $20$
via family $\eta$ is answered without computation: $\eta$'s discriminant is
$-16$. Further instances therefore require \emph{non-sporadic} rectified
families at genus-zero levels with larger cusp spaces --- the reverse
factory's next search domain, now equipped with a quantitative objective
(maximize the error-decay base against the denominator-growth base).
\end{remark}

\subsection{The general conjecture}
\begin{conjecture}[D-2, in Sol's refined mechanism form]\label{conj:d2}
For a modular Picard--Fuchs family with a fold fixed by an Atkin--Lehner
involution, a rational integral source in the $(-1)$-eigenspace that
vanishes to the required order at the fold yields a regularized companion
whose connection value lies in the odd critical-period line. Quantitatively:
for every genus-zero level $N$ with $\dim S_{r+1}\ge2$ and a rectified
family whose fold is the $W_N$-fixed point, the Fricke-odd vanishing
combination yields a $d_n^r$-integral, exponentially convergent apparatus
whose limit is a rational multiple of an odd critical $L$-value.
\end{conjecture}

Status: \Open. Evidence: two data points (levels $6$ and $12$, the second
positive and the first explaining the necessity of the parity hypothesis),
plus the $\delta$ control isolating the vanishing clause, plus the three
separated layers --- \emph{integral structure} $\Rightarrow d_n^r$;
\emph{fold-vanishing} $\Rightarrow$ regularization; \emph{odd parity}
$\Rightarrow$ odd critical value --- each of which is independently
supported. Two data points is not evidence for a general law, and we do
not present it as such; the value of the conjecture is that it states
exactly which hypotheses a third data point must test.

%======================================================================
\section{Honest irrationality bookkeeping}\label{sec:race}

It would be easy to over-read \S\ref{sec:level12}. We therefore state the
negative plainly.

\subsection{The apparatus loses the race}
An Ap\'ery-style irrationality proof needs the cleared error to tend to
zero:
\[
\bigl|\,d_n^{3}\bigl(A_n\,\xi^*-B^*_n\bigr)\,\bigr|\;\longrightarrow\;0 .
\]
Here $d_n^3\sim e^{3n}=(20.08\ldots)^n$ while $A_n|B^*_n/A_n-\xi^*|$ decays
only like $4^{-n}$ relative to $A_n$; the product behaves like
$e^{3n}\cdot4^{-n}\cdot(\text{subexponential})\to\infty$.

\begin{observation}[no irrationality]\label{obs:noirr}
\Proved{} The apparatus of Construction~\ref{con:app} does \emph{not}
prove, and does not bear on, the irrationality of $L(f_6,3)$. The
decay base $4$ loses to the denominator base $e^3=20.08\ldots$.
\end{observation}

Its value is structural: cuspidal critical $L$-values are now inside the
modular-anchor factory with Ap\'ery-quality arithmetic, and the
convergence/denominator trade-off becomes an \emph{optimization problem}
over the triple (level, family, source) --- the same problem that Ap\'ery's
$\zeta(3)$ construction wins in the Eisenstein class.

\subsection{The race, scored over all fifteen families}
Sol's scoring criterion: with $A_n\sim\alpha^n$, error $\sim\beta^n$, and
denominators $\ll\delta^n$, one needs $\delta|\beta|<1$. For our families
$|A_n\xi-B_n|\sim|t_c'|^{-n}$ and $\delta=e^{r}$, so the score is
$e^{r}/|t_c'|$ with $t_c'$ the larger root of $P$.

\begin{verification}[Route-1 scan]\label{ver:race}
\VerifiedR{exact roots, all fifteen families}
\begin{center}
\begin{tabular}{@{}llllll@{}}
\toprule
family & $r$ & $|t_c'|$ & $e^{r}/|t_c'|$ & verdict \\
\midrule
$D$ ($\zeta(2)$)      & 2 & $11.09$ & $0.67$ & \textbf{win} --- Ap\'ery's $\zeta(2)$ \\
$\gamma$ ($\zeta(3)$) & 3 & $33.97$ & $0.59$ & \textbf{win} --- Ap\'ery's $\zeta(3)$ \\
$\varepsilon$         & 3 & $1.457$ & $13.8$ & fail (best of the rest) \\
all others            &   & $\le1$  & $\ge7.4$ & fail \\
\bottomrule
\end{tabular}
\end{center}
\end{verification}

Exactly the two classical wins, with nothing near the boundary. The
sporadic table is exhausted as a source of irrationality proofs.

\subsection{Ap\'ery's miracle is the same construction}
\begin{verification}[the classical source vanishes at its fold]\label{ver:phigamma}
\VerifiedR{$|\Phi_\gamma(\tau_c)|<10^{-122}$}
$\Phi_\gamma(\tau_c)=0$: Ap\'ery's own $\zeta(3)$ source vanishes at the
fold. It is the $W_6$-\emph{odd} Eisenstein direction, and the vanishing is
Corollary~\ref{cor:oddvanish} again --- any weight-$4$ form odd under
$W_6$ vanishes at $i/\sqrt6$ since the automorphy factor there is $+1$.
\end{verification}

This reframes the whole picture. Ap\'ery's construction \emph{is} the
construction of \S\ref{sec:level12}: an odd source vanishing at an
Atkin--Lehner-fixed fold. It succeeds where ours does not only because its
level has root ratio $34>e^3$. One mechanism, two classes:

\begin{center}
\begin{tabular}{@{}lll@{}}
\toprule
& odd Eisenstein source & odd cusp source \\
\midrule
level $6$ (ratio $34>e^3$) & Ap\'ery's $\zeta(3)$ proof &
$S_4^-(\Gamma_0(6))=0$ --- none exists \\
level $12$ (ratio $4<e^3$) & $\zeta(3)$-type, loses race &
our $L(f_6,3)/2$ apparatus, loses race \\
\bottomrule
\end{tabular}
\end{center}

Note also that the ``mixed-source fix'' at level $6$ --- subtract a
multiple of the Eisenstein source to force $f_6$ to vanish at the fold ---
is impossible: $\Phi_\gamma$ \emph{already} vanishes there, so no multiple
of it changes the value of $f_6$ at $\tau_c$ \ExcludedR{verified}.

\subsection{The precise open problem}
\begin{problem}[the cuspidal-Ap\'ery question]\label{prob:cusp}
\Open{} Find a genus-zero group $\Gamma$ carrying a rectified family whose
Atkin--Lehner-fixed fold has conjugate root $|t_c'|>e^{r}$, and with
$S^-_{r+1}(\Gamma)\ne0$. The odd cuspidal companion would then give a
race-winning $d_n^r$-apparatus for a cusp-form critical $L$-value --- an
irrationality proof of a cuspidal period --- \emph{provided} the
$d_n^r$-integrality, which is proved in the Eisenstein class but only
observed in the cuspidal one, survives.
\end{problem}

Two failure modes deserve respect. First, cuspidal modular-symbol
denominators may be genuinely worse than $d_n^r$ (Sol's warning); our
$n\le40$ data cannot see this. Second, the sporadic table itself is
exhausted (Verification~\ref{ver:race}, Corollary~\ref{cor:unique}), so the
search must move to non-sporadic rectified families.

%======================================================================
\section{Evidence status}\label{sec:status}

Consolidated, with nothing upgraded relative to the ledgers.

\begin{enumerate}[leftmargin=2.4em,label=(\arabic*)]
\item \textbf{The Eisenstein/cuspidal dichotomy of sources}
(\S\ref{sec:intro}): \VerifiedR{exact $\Q$, coefficientwise to $q^{60}$,
held-out band $q^{27}$--$q^{60}$}, conditional on the modularity
certification of $t,F$ per family (complete only for level $9$).
\item \textbf{Level-$6$ identity} $f_6\sqrt{1-34t+t^2}=\Phi_\gamma$:
\VerifiedR{$q^{60}$}, theorem-grade modulo the same certification, via the
Sturm bound $8$ in $M_8(\Gamma_0(6))$.
\item \textbf{Level-$6$ limit} $\Theta(\tau_c)=L(f_6,3)-\tfrac\pi{\sqrt6}
L(f_6,2)$: \VerifiedR{$120$ digits}; period-polynomial proof route stated,
not written out. Second-kind constant $\Theta_2(\tau_c)$: \Open{} (PSLQ
negative over the stated bases).
\item \textbf{Oldspace swap and $\fs(\tau_{12})=0$}: \ProvedS{}
(Theorem~\ref{thm:oldspace}); previously \VerifiedR{PSLQ $35^+$ digits}.
\item \textbf{Source identity} $\fs\sqrt{1-4t}=\Phi_\alpha$:
\VerifiedR{exact to $q^{40}$} via the central-binomial reconstruction of
$R^*$; Sturm-boundable in $M_8(\Gamma_0(12))$ (bound $16$), which the
verification range exceeds.
\item \textbf{$d_n^3$-integrality} of $B^f$ and $B^*$: \VerifiedR{$n\le40$}
in both cases; \Open{} beyond. No denominator theorem is claimed.
\item \textbf{The limit} $B^*_n/A_n=L(f_6,3)/2$: \VerifiedR{$61$-digit
PSLQ, residual $2\times10^{-61}$}, with $121$-digit stability of the ratio;
the shape (rational multiple of the odd critical value) is supported by
Propositions~\ref{prop:parity}--\ref{prop:half}, the factor $\tfrac12$ is
\RouteOpen.
\item \textbf{Mechanism}: Effect 1 \Proved{} given item (5)
(Lemma~\ref{lem:regular}); Effect 2 structurally \ProvedS{}
(Proposition~\ref{prop:parity}), exact constant \RouteOpen.
\item \textbf{$\delta$ control}: \VerifiedR{$n\le40$ exact, $120$-digit
numerics}. \textbf{Uniqueness of $\alpha$}: \Proved{} given the tabulated
fold and dimension data.
\item \textbf{$\kappa$-cancellation}: \ExcludedR{PSLQ, $40$ digits, stated
bases and height caps}.
\item \textbf{Race scan and} $\Phi_\gamma(\tau_c)=0$:
\VerifiedR{exact roots; $10^{-122}$}.
\item \textbf{Conjecture D-2}: \Open, two data points.
\item \textbf{Irrationality}: none. Observation~\ref{obs:noirr} is a
proved negative about \emph{this} apparatus, not about $L(f_6,3)$.
\end{enumerate}

\begin{remark}[for the newcomer]
The pattern to internalize: a $61$-digit PSLQ identification is
overwhelming \emph{evidence} and zero \emph{proof}; a $q^{40}$ coefficient
agreement inside a space with Sturm bound $16$ \emph{is} a proof, once the
ambient space is certified; and an integrality check to $n=40$ is neither,
because denominator theorems in this subject fail first at moderate $n$ for
reasons invisible in small data. Keeping these three apart is the whole of
the discipline.
\end{remark}

%======================================================================
\section{Provenance and data}\label{sec:prov}

All computations are logged in the programme's ledgers, with exact
arithmetic (Python \texttt{Fraction}) for series work and \texttt{mpmath}
at $120$ decimal digits for the numerics.

\begin{sloppypar}
\begin{itemize}[leftmargin=1.6em]
\item Source dichotomy and the twelve identifications:
\texttt{work/PHI\_SOURCE\_LEDGER.md}; scripts
\texttt{work/z5eps/eps60\_phi\_source.py},
\texttt{eps60b\_phi\_holdouts.py}, \texttt{eps61\_eichler\_limits.py}.
\item Level-$6$ experiment (\S\ref{sec:level6}):
\texttt{work/CUSPIDAL\_COMPANION.md}; scripts
\texttt{work/z5eps/eps62\_cuspidal\_companion.py},
\texttt{eps62b\_cusp\_limit.py}.
\item Level-$12$ construction, $\delta$ control, uniqueness, proof notes,
race scan (\S\ref{sec:level12}--\S\ref{sec:race}):
\texttt{work/D1\_CUSPIDAL\_APPARATUS.md}, \S\S1--9.
\item Proved base (rectification, companion inversion, uniform companion
formula): \texttt{papers\_out/modular\_anchors/main.tex}.
\end{itemize}
\end{sloppypar}

\subsection*{Acknowledgements}
The collaborator Sol proposed the cuspidal-source project, formulated open
problem D-1 and the six-item theorem package, supplied the refined
hypotheses of Conjecture~\ref{conj:d2}, directed the $\delta$ control, and
devised the irrationality-race scoring used in \S\ref{sec:race}; the
reviews were adversarial in the productive sense and several negative
results in this paper exist because of them.

\begin{thebibliography}{9}
\bibitem{anchors} \emph{Modular anchors for Ap\'ery-like companions:
second solutions as Eichler integrals, with the symmetric-square theorem
for the sporadic operators}, companion paper of this programme,
\texttt{papers\_out/modular\_anchors}.
\bibitem{philedger} \emph{The companion sources $\Phi$ of the fifteen
sporadic pairs}, programme ledger \texttt{work/PHI\_SOURCE\_LEDGER.md}.
\bibitem{apery} R.~Ap\'ery, \emph{Irrationalit\'e de $\zeta(2)$ et
$\zeta(3)$}, Ast\'erisque \textbf{61} (1979), 11--13.
\bibitem{beukers} F.~Beukers, \emph{Irrationality proofs using modular
forms}, Ast\'erisque \textbf{147--148} (1987), 271--283.
\bibitem{zagier} D.~Zagier, \emph{Integral solutions of Ap\'ery-like
recurrence equations}, in: Groups and Symmetries, CRM Proc.\ Lecture Notes
\textbf{47} (2009), 349--366.
\bibitem{alw} A.~O.~L.~Atkin, J.~Lehner, \emph{Hecke operators on
$\Gamma_0(m)$}, Math.\ Ann.\ \textbf{185} (1970), 134--160.
\bibitem{sturm} J.~Sturm, \emph{On the congruence of modular forms},
Lecture Notes in Math.\ \textbf{1240}, Springer (1987), 275--280.
\bibitem{cooper} S.~Cooper, \emph{Sporadic sequences, modular forms and
new series for $1/\pi$}, Ramanujan J.\ \textbf{29} (2012), 163--183.
\end{thebibliography}

\end{document}
